% !TeX program = XeLaTeX
% !TeX root = ../pujavidhanam.tex
\chapt{माघ-स्नान-विधिः}

\centerline{\small{(मूलम्—श्रीव्रतरत्नाकरः, प्रथमो भागः)}}

\sect{स्नानम्}

\twolineshloka
{माघमासे रटन्त्यापः किञ्चिदभ्युदिते रवौ}
{ब्रह्मघ्नं वा सुरापं वा यं पतन्तं पुनीमहे}

\twolineshloka
{दुःखदारिद्र्यनाशाय श्रीविष्णोस्तोषणाय च}
{प्रातः स्नानं करोम्यद्य माघे पापविनाशनम्}

\twolineshloka
{मकरस्थे रवौ माघे गोविन्दाच्युत माधव}
{स्नानेनानेन मे देव यथोक्तफलदो भव}

\twolineshloka
{कृष्णार्जुन निमज्जामि प्रभातेऽस्मिन् शुभोदके}
{अनेन माघस्नानेन सुप्रीतो मां समुद्धर}

\dnsub{प्रत्यहं सूर्यार्घ्यमन्त्रः}

\twolineshloka
{सवित्रे प्रसवित्रे च परं धाम जले मम}
{त्वत्तेजसा परिभ्रष्टं पापं यातु सहस्रधा}
\nobreak\hfill{}—सवित्रे नमः इदमर्घ्यम्। (त्रिः)

\dnsub{त्रिवेण्यर्घ्यम्}

\twolineshloka
{गङ्गायमुनयोर्मध्ये यत्र गुप्ता सरस्वती}
{त्रैलोक्यवन्दिते देवि त्रिवेण्यर्घ्यं ददामि ते}
\nobreak\hfill{}—त्रिवेण्यै नमः इदमर्घ्यम्। (त्रिः)

\dnsub{प्रार्थना}

\twolineshloka*
{प्रभाकर जगन्नाथ दिवाकर नमोऽस्तु ते}
{परिपूर्णं कुरुष्वेदं माघस्नानं मया कृतम्}

इति सूर्यं प्रार्थयेत्॥

\closesection
